% From sections/sec07_five.tex, cut 2026-10-19 (Remark rem:interior; Theorem thm:refute with part (b) and the first-exceeds paragraph; Neuwirth sentence and Lemma lem:tenvertices; old Theorem thm:fivefive with the connected-one-cycle claim).
\begin{remark}[the padding row has to be a boundary row]\label{rem:interior}
Take $h=3$, $B_1=B_2=\{1\}$, $B_3=\{1,3\}$, and padding row $2$. Columns $0$ and $-1$
keep the board valid, but $-2$ breaks it: the diamond $(-2,2)$,
$(-1,1)$, $(0,2)$, $(-1,3)$ and the rectangle on columns $-1$, $3$ and rows $1$, $3$ both
show $\{(-1,1),(-1,3)\}$.
\end{remark}

\begin{theorem}[the guess fails, computer-assisted]\label{thm:refute}
\begin{enumerate}[(a)]
\item We have $f(6,6)=12$, $f(6,7)=13$, $f(7,7)=14$, and $f(6,8)=15$. The blackout $(1,1)$,
$(1,5)$, $(2,1)$, $(3,1)$, $(3,2)$, $(3,3)$, $(4,3)$, $(4,4)$, $(5,4)$, $(6,4)$, $(7,4)$,
$(7,5)$, $(7,6)$, $(8,2)$, $(8,6)$ of the $6\times8$ grid attains the last value.
\item The guess of Section~\ref{sec:intro} fails. Indeed $f(5,8)=14$ and $f(6,8)=15$.
\end{enumerate}
\end{theorem}

\begin{proof}
(a) The other lower bounds are stored blackouts of $h+m$ points. Two SAT encodings find no
valid blackout of $13$, $14$, $15$, or $16$ points on the $6\times6$, $6\times7$,
$7\times7$, or $6\times8$ grid, respectively, and all four bounds have checked certificates
(Section~\ref{sec:verify}).

(b) The blackouts $S_8$ and that of part~(a) have $14$ and $15$ points, against the guessed
$13$ and $14$. Checking their $364$ and $592$ rectangles confirms validity, so the
refutation uses no solver.
\end{proof}

If a side is at most four, then $f(h,m)\le h+m$ by Sections~\ref{sec:strip}
to~\ref{sec:four}. Also $f(h,m)=h+m$ on $5\times5$, $5\times6$, $5\times7$, and $6\times6$,
the other grids of area below $40$. So $f$ first exceeds $h+m$ on the $5\times8$ grid.

Neuwirth~\cite[Remark~18]{neuwirth2001} states the case $5+5$ of the next lemma without
proof.

\begin{lemma}[few edges without short cycles]\label{lem:tenvertices}
Let $G$ be a finite simple bipartite graph with $v$ vertices, $1\le v\le10$, and with no cycle
of length four or six. Then $G$ has at most $v$ edges. If it has exactly $v$ edges, then $G$ is
connected and has exactly one cycle, of length $8$ or $10$.
\end{lemma}

\begin{proof}
Every cycle of $G$ has length at least $8$. Suppose a component has more edges than
vertices. Removing vertices of degree one with their edges until none is left gives a connected
$H$ with more edges than vertices and minimum degree two, so some vertex $x$ of $H$ has degree
at least three. Let $D_k$ be the set of vertices at distance $k$ from $x$. Two paths of length
at most $3$ from $x$ to one vertex would close a cycle of length at most $6$. So for $k=1,2,3$
every vertex of $D_k$ has exactly one neighbor in $D_{k-1}$ and none in $D_k$. Its degree is
at least two, so it also has a neighbor in $D_{k+1}$. For $k\le2$ no other vertex of $D_k$
shares that neighbor. Hence $|D_1|\ge3$, $|D_3|\ge|D_2|\ge|D_1|$ and $|D_4|\ge1$, at least
$11$ vertices in all. So $G$ has at most $v$ edges, and with exactly $v$ edges every
component has one cycle. Two disjoint cycles of length at least $8$ need $16$ vertices, so $G$ is connected.
\end{proof}

\begin{theorem}[the $5\times5$ grid]\label{thm:fivefive}
\begin{enumerate}[(a)]
\item We have $f(5,5)=10$, and $G_S$ is connected with exactly one cycle for every optimum $S$.
\item (Computer-assisted.) There are exactly $a_5(5)=8$ maximum blackouts, the eight distinct
images of $S_0$ below under the symmetries of the square.
\end{enumerate}
\end{theorem}

\begin{proof}
\emph{Upper bound.} If $S$ is valid, $G_S$ has $10$ vertices, $|S|$ edges and no cycle of
length four or six (Lemma~\ref{lem:girth}), so Lemma~\ref{lem:tenvertices} applies.

\emph{Lower bound.} Let $S_0$ have $B_1=\{4,5\}$, $B_2=\{3,4\}$, $B_3=\{1,2,3\}$, $B_4=\{1\}$, and $B_5=\{1,5\}$. Its graph is an $8$-cycle with two pendant
edges, so $S_0$ is axis-valid. By Lemma~\ref{lem:leftmost} the $30$ tilted rectangles among
the $130$ are the nine unit diamonds, the four kinds of Lemma~\ref{lem:tilted-narrow} in four
positions each, and five that span the whole grid, one of them a large diamond. Exactly fifteen are weak. The
table lists the two corners each one shows and, after a bar for an oblique pair, the white
corner of its axis-aligned rectangle.
\begin{center}\footnotesize
\begin{tabular}{@{}llll@{}}
$(1,2),(2,1)\mid(1,1)$ & $(1,3),(2,2)\mid(1,2)$ & $(3,4),(4,3)\mid(4,4)$ & $(2,5),(3,4)\mid(3,5)$ \\
$(2,2),(4,2)$ & $(1,2),(4,3)\mid(1,3)$ & $(4,4),(5,2)\mid(4,2)$ & $(3,5),(4,4)\mid(3,4)$ \\
$(4,3),(5,2)\mid(4,2)$ & $(3,4),(5,2)\mid(5,4)$ & $(2,5),(5,2)\mid(2,2)$ & $(4,5),(5,3)\mid(4,3)$ \\
$(1,3),(4,2)\mid(1,2)$ & $(2,5),(4,3)\mid(4,5)$ & & $(4,5),(5,4)\mid(4,4)$ \\
\end{tabular}
\end{center}
The fifteen pairs are distinct, and the other tilted rectangles are determined by
Lemma~\ref{lem:threecorners}. Each listed white corner lies on the only axis-aligned rectangle
through its pair, and $(2,2),(4,2)$ has no rival since $B_2\cap B_4=\emptyset$. So $S_0$ is valid.

\emph{(b).} Of the $44{,}640$ sets of $10$ points without a short cycle, exactly the $8$
images of $S_0$ are valid, by a direct check confirmed by the ILP program and brute force.
\end{proof}

% From sections/sec07_five.tex, 2026-10-19: pre-compression text of the collar proof, the five-row theorem and the six-row proof, and the old figure placement.
\section{Five and six rows}\label{sec:five}

\subsection{A finite test for boundary padding}

For a blackout $S$ of the $h\times m$ grid and $L,R\ge0$, let $S_{L,R}$ add the columns
$1-L,\dots,0$ and $m+1,\dots,m+R$, each black exactly in the \emph{padding row}, row $1$ or
row $h$. Here $W_h$ is the span bound of Lemma~\ref{lem:span}.

\begin{theorem}[finite collar]\label{thm:collar}
Let $h\ge3$ and $m\ge1$, and let $S$ be a blackout of the $h\times m$ grid. Fix the padding
row and choose the left side, the right side, or both. Call $(L,R)$ \emph{admissible} if
$L$ and $R$ vanish on the sides not chosen, and let $L^*$ and $R^*$ be $W_h$ on the chosen
sides and $0$ on the others. Then $S_{L,R}$ is valid for every admissible $(L,R)$ if and only if $S_{L^*,R^*}$ is
valid.
\end{theorem}

\begin{proof}
Write $S^*=S_{L^*,R^*}$. If every admissible padding is valid, then so is $S^*$.
Conversely, let $S^*$ be valid and
$U=S_{L,R}$ with $(L,R)$ admissible. By Lemma~\ref{lem:leftmost} the corners of a tilted
rectangle lie in the rows $j$, $j-b$, $j+d$, $j+d-b$ with $b,d\ge1$, so at most one lies in
row $1$ and at most one in row $h$. Hence a weak tilted rectangle of $U$ has a black corner in an
original column $e$, and by Lemma~\ref{lem:span} its columns lie in
$[e-W_h,e+W_h]\cap[1-L,m+R]\subseteq[1-L^*,m+R^*]$. Since $U$ and $S^*$ agree on their common
columns, it is a rectangle of $S^*$ with the same presentation. Suppose that two distinct
rectangles $R_1$ and $R_2$ of $U$ have the same presentation $P$.

\emph{Step 1. Both axis-aligned.} The graph $G_U$ is $G_S$ plus one vertex of degree one for
each added column, so its cycles are cycles of $G_{S^*}$. Since $S^*$ is valid with at least
three columns, neither graph has a cycle of length four or six, and $U$ is axis-valid
(Lemma~\ref{lem:girth}).

\emph{Step 2. Both tilted.} By Lemma~\ref{lem:threecorners}, $|P|\le2$, so both are weak and
are rectangles of $S^*$ with the same presentation, which is impossible.

\emph{Step 3. One of each.} Let $R_1$ be tilted, so it is a rectangle of $S^*$. A column of
$R_2$ with a point of $P$ is a column of $R_1$, and one without has two black corners and is
an original column. So $R_2$ is a rectangle of $S^*$ too, a contradiction.
\end{proof}
\subsection{Five rows}

Let $S_8$ be the blackout of the $5\times8$ grid with $B_1=\{1,4\}$, $B_2=\{4,5\}$,
$B_3=\{2,5\}$, $B_4=B_5=\{5\}$, $B_6=\{3,4\}$, $B_7=\{2,3\}$, and $B_8=\{1,2\}$. Its
doubles carry the edges of $K_{2,3}$ between $\{2,4\}$ and $\{1,3,5\}$.

\begin{figure}[htbp]
\centering
\begin{tikzpicture}[x=0.8cm,y=0.8cm,baseline=(current bounding box.north)]
  \begin{scope}[xshift=0.000cm]
    \draw[black,line width=0.32pt,dash pattern=on 1.20pt off 1.20pt] (1,-1) -- (13,-1) (1,-2) -- (13,-2) (1,-3) -- (13,-3) (1,-4) -- (13,-4) (1,-5) -- (13,-5) (1,-1) -- (1,-5) (2,-1) -- (2,-5) (3,-1) -- (3,-5) (4,-1) -- (4,-5) (5,-1) -- (5,-5) (6,-1) -- (6,-5) (7,-1) -- (7,-5) (8,-1) -- (8,-5) (9,-1) -- (9,-5) (10,-1) -- (10,-5) (11,-1) -- (11,-5) (12,-1) -- (12,-5) (13,-1) -- (13,-5);
    \filldraw[fill=kgray,draw=black,line width=0.48pt] (1,-1) circle (0.3);
    \draw[black,line width=0.48pt] (0.788,-0.788) -- (1.212,-1.212) (0.788,-1.212) -- (1.212,-0.788);
    \filldraw[fill=white,draw=black,line width=0.48pt] (2,-1) circle (0.3);
    \filldraw[fill=white,draw=black,line width=0.48pt] (3,-1) circle (0.3);
    \filldraw[fill=white,draw=black,line width=0.48pt] (4,-1) circle (0.3);
    \filldraw[fill=white,draw=black,line width=0.48pt] (5,-1) circle (0.3);
    \filldraw[fill=white,draw=black,line width=0.48pt] (6,-1) circle (0.3);
    \filldraw[fill=white,draw=black,line width=0.48pt] (7,-1) circle (0.3);
    \filldraw[fill=kgray,draw=black,line width=0.48pt] (8,-1) circle (0.3);
    \draw[black,line width=0.48pt] (7.788,-0.788) -- (8.212,-1.212) (7.788,-1.212) -- (8.212,-0.788);
    \filldraw[fill=kgray,draw=black,line width=0.48pt] (9,-1) circle (0.3);
    \draw[black,line width=0.48pt] (8.788,-0.788) -- (9.212,-1.212) (8.788,-1.212) -- (9.212,-0.788);
    \filldraw[fill=kgray,draw=black,line width=0.48pt] (10,-1) circle (0.3);
    \draw[black,line width=0.48pt] (9.788,-0.788) -- (10.212,-1.212) (9.788,-1.212) -- (10.212,-0.788);
    \filldraw[fill=kgray,draw=black,line width=0.48pt] (11,-1) circle (0.3);
    \draw[black,line width=0.48pt] (10.788,-0.788) -- (11.212,-1.212) (10.788,-1.212) -- (11.212,-0.788);
    \filldraw[fill=kgray,draw=black,line width=0.48pt] (12,-1) circle (0.3);
    \draw[black,line width=0.48pt] (11.788,-0.788) -- (12.212,-1.212) (11.788,-1.212) -- (12.212,-0.788);
    \filldraw[fill=kgray,draw=black,line width=0.48pt] (13,-1) circle (0.3);
    \draw[black,line width=0.48pt] (12.788,-0.788) -- (13.212,-1.212) (12.788,-1.212) -- (13.212,-0.788);
    \filldraw[fill=white,draw=black,line width=0.48pt] (1,-2) circle (0.3);
    \filldraw[fill=white,draw=black,line width=0.48pt] (2,-2) circle (0.3);
    \filldraw[fill=kgray,draw=black,line width=0.48pt] (3,-2) circle (0.3);
    \draw[black,line width=0.48pt] (2.788,-1.788) -- (3.212,-2.212) (2.788,-2.212) -- (3.212,-1.788);
    \filldraw[fill=white,draw=black,line width=0.48pt] (4,-2) circle (0.3);
    \filldraw[fill=white,draw=black,line width=0.48pt] (5,-2) circle (0.3);
    \filldraw[fill=white,draw=black,line width=0.48pt] (6,-2) circle (0.3);
    \filldraw[fill=kgray,draw=black,line width=0.48pt] (7,-2) circle (0.3);
    \draw[black,line width=0.48pt] (6.788,-1.788) -- (7.212,-2.212) (6.788,-2.212) -- (7.212,-1.788);
    \filldraw[fill=kgray,draw=black,line width=0.48pt] (8,-2) circle (0.3);
    \draw[black,line width=0.48pt] (7.788,-1.788) -- (8.212,-2.212) (7.788,-2.212) -- (8.212,-1.788);
    \filldraw[fill=white,draw=black,line width=0.48pt] (9,-2) circle (0.3);
    \filldraw[fill=white,draw=black,line width=0.48pt] (10,-2) circle (0.3);
    \filldraw[fill=white,draw=black,line width=0.48pt] (11,-2) circle (0.3);
    \filldraw[fill=white,draw=black,line width=0.48pt] (12,-2) circle (0.3);
    \filldraw[fill=white,draw=black,line width=0.48pt] (13,-2) circle (0.3);
    \filldraw[fill=white,draw=black,line width=0.48pt] (1,-3) circle (0.3);
    \filldraw[fill=white,draw=black,line width=0.48pt] (2,-3) circle (0.3);
    \filldraw[fill=white,draw=black,line width=0.48pt] (3,-3) circle (0.3);
    \filldraw[fill=white,draw=black,line width=0.48pt] (4,-3) circle (0.3);
    \filldraw[fill=white,draw=black,line width=0.48pt] (5,-3) circle (0.3);
    \filldraw[fill=kgray,draw=black,line width=0.48pt] (6,-3) circle (0.3);
    \draw[black,line width=0.48pt] (5.788,-2.788) -- (6.212,-3.212) (5.788,-3.212) -- (6.212,-2.788);
    \filldraw[fill=kgray,draw=black,line width=0.48pt] (7,-3) circle (0.3);
    \draw[black,line width=0.48pt] (6.788,-2.788) -- (7.212,-3.212) (6.788,-3.212) -- (7.212,-2.788);
    \filldraw[fill=white,draw=black,line width=0.48pt] (8,-3) circle (0.3);
    \filldraw[fill=white,draw=black,line width=0.48pt] (9,-3) circle (0.3);
    \filldraw[fill=white,draw=black,line width=0.48pt] (10,-3) circle (0.3);
    \filldraw[fill=white,draw=black,line width=0.48pt] (11,-3) circle (0.3);
    \filldraw[fill=white,draw=black,line width=0.48pt] (12,-3) circle (0.3);
    \filldraw[fill=white,draw=black,line width=0.48pt] (13,-3) circle (0.3);
    \filldraw[fill=kgray,draw=black,line width=0.48pt] (1,-4) circle (0.3);
    \draw[black,line width=0.48pt] (0.788,-3.788) -- (1.212,-4.212) (0.788,-4.212) -- (1.212,-3.788);
    \filldraw[fill=kgray,draw=black,line width=0.48pt] (2,-4) circle (0.3);
    \draw[black,line width=0.48pt] (1.788,-3.788) -- (2.212,-4.212) (1.788,-4.212) -- (2.212,-3.788);
    \filldraw[fill=white,draw=black,line width=0.48pt] (3,-4) circle (0.3);
    \filldraw[fill=white,draw=black,line width=0.48pt] (4,-4) circle (0.3);
    \filldraw[fill=white,draw=black,line width=0.48pt] (5,-4) circle (0.3);
    \filldraw[fill=kgray,draw=black,line width=0.48pt] (6,-4) circle (0.3);
    \draw[black,line width=0.48pt] (5.788,-3.788) -- (6.212,-4.212) (5.788,-4.212) -- (6.212,-3.788);
    \filldraw[fill=white,draw=black,line width=0.48pt] (7,-4) circle (0.3);
    \filldraw[fill=white,draw=black,line width=0.48pt] (8,-4) circle (0.3);
    \filldraw[fill=white,draw=black,line width=0.48pt] (9,-4) circle (0.3);
    \filldraw[fill=white,draw=black,line width=0.48pt] (10,-4) circle (0.3);
    \filldraw[fill=white,draw=black,line width=0.48pt] (11,-4) circle (0.3);
    \filldraw[fill=white,draw=black,line width=0.48pt] (12,-4) circle (0.3);
    \filldraw[fill=white,draw=black,line width=0.48pt] (13,-4) circle (0.3);
    \filldraw[fill=white,draw=black,line width=0.48pt] (1,-5) circle (0.3);
    \filldraw[fill=kgray,draw=black,line width=0.48pt] (2,-5) circle (0.3);
    \draw[black,line width=0.48pt] (1.788,-4.788) -- (2.212,-5.212) (1.788,-5.212) -- (2.212,-4.788);
    \filldraw[fill=kgray,draw=black,line width=0.48pt] (3,-5) circle (0.3);
    \draw[black,line width=0.48pt] (2.788,-4.788) -- (3.212,-5.212) (2.788,-5.212) -- (3.212,-4.788);
    \filldraw[fill=kgray,draw=black,line width=0.48pt] (4,-5) circle (0.3);
    \draw[black,line width=0.48pt] (3.788,-4.788) -- (4.212,-5.212) (3.788,-5.212) -- (4.212,-4.788);
    \filldraw[fill=kgray,draw=black,line width=0.48pt] (5,-5) circle (0.3);
    \draw[black,line width=0.48pt] (4.788,-4.788) -- (5.212,-5.212) (4.788,-5.212) -- (5.212,-4.788);
    \filldraw[fill=white,draw=black,line width=0.48pt] (6,-5) circle (0.3);
    \filldraw[fill=white,draw=black,line width=0.48pt] (7,-5) circle (0.3);
    \filldraw[fill=white,draw=black,line width=0.48pt] (8,-5) circle (0.3);
    \filldraw[fill=white,draw=black,line width=0.48pt] (9,-5) circle (0.3);
    \filldraw[fill=white,draw=black,line width=0.48pt] (10,-5) circle (0.3);
    \filldraw[fill=white,draw=black,line width=0.48pt] (11,-5) circle (0.3);
    \filldraw[fill=white,draw=black,line width=0.48pt] (12,-5) circle (0.3);
    \filldraw[fill=white,draw=black,line width=0.48pt] (13,-5) circle (0.3);
    \node[lab] at (1,-0.45) {1};
    \node[lab] at (2,-0.45) {2};
    \node[lab] at (3,-0.45) {3};
    \node[lab] at (4,-0.45) {4};
    \node[lab] at (5,-0.45) {5};
    \node[lab] at (6,-0.45) {6};
    \node[lab] at (7,-0.45) {7};
    \node[lab] at (8,-0.45) {8};
    \node[lab] at (9,-0.45) {9};
    \node[lab] at (10,-0.45) {10};
    \node[lab] at (11,-0.45) {11};
    \node[lab] at (12,-0.45) {12};
    \node[lab] at (13,-0.45) {13};
    \node[lab] at (0.45,-1) {1};
    \node[lab] at (0.45,-2) {2};
    \node[lab] at (0.45,-3) {3};
    \node[lab] at (0.45,-4) {4};
    \node[lab] at (0.45,-5) {5};
    \node[capt,anchor=north] at (7.00,-5.45) {(a)};
  \end{scope}
\end{tikzpicture}

\par\medskip
\begin{tikzpicture}[x=0.8cm,y=0.8cm,baseline=(current bounding box.north)]
  \begin{scope}[xshift=0.000cm]
    \draw[black,line width=0.32pt,dash pattern=on 1.20pt off 1.20pt] (1,-1) -- (16,-1) (1,-2) -- (16,-2) (1,-3) -- (16,-3) (1,-4) -- (16,-4) (1,-5) -- (16,-5) (1,-6) -- (16,-6) (1,-1) -- (1,-6) (2,-1) -- (2,-6) (3,-1) -- (3,-6) (4,-1) -- (4,-6) (5,-1) -- (5,-6) (6,-1) -- (6,-6) (7,-1) -- (7,-6) (8,-1) -- (8,-6) (9,-1) -- (9,-6) (10,-1) -- (10,-6) (11,-1) -- (11,-6) (12,-1) -- (12,-6) (13,-1) -- (13,-6) (14,-1) -- (14,-6) (15,-1) -- (15,-6) (16,-1) -- (16,-6);
    \filldraw[fill=white,draw=black,line width=0.48pt] (1,-1) circle (0.3);
    \filldraw[fill=kgray,draw=black,line width=0.48pt] (2,-1) circle (0.3);
    \draw[black,line width=0.48pt] (1.788,-0.788) -- (2.212,-1.212) (1.788,-1.212) -- (2.212,-0.788);
    \filldraw[fill=white,draw=black,line width=0.48pt] (3,-1) circle (0.3);
    \filldraw[fill=white,draw=black,line width=0.48pt] (4,-1) circle (0.3);
    \filldraw[fill=white,draw=black,line width=0.48pt] (5,-1) circle (0.3);
    \filldraw[fill=kgray,draw=black,line width=0.48pt] (6,-1) circle (0.3);
    \draw[black,line width=0.48pt] (5.788,-0.788) -- (6.212,-1.212) (5.788,-1.212) -- (6.212,-0.788);
    \filldraw[fill=white,draw=black,line width=0.48pt] (7,-1) circle (0.3);
    \filldraw[fill=white,draw=black,line width=0.48pt] (8,-1) circle (0.3);
    \filldraw[fill=white,draw=black,line width=0.48pt] (9,-1) circle (0.3);
    \filldraw[fill=white,draw=black,line width=0.48pt] (10,-1) circle (0.3);
    \filldraw[fill=white,draw=black,line width=0.48pt] (11,-1) circle (0.3);
    \filldraw[fill=white,draw=black,line width=0.48pt] (12,-1) circle (0.3);
    \filldraw[fill=white,draw=black,line width=0.48pt] (13,-1) circle (0.3);
    \filldraw[fill=white,draw=black,line width=0.48pt] (14,-1) circle (0.3);
    \filldraw[fill=white,draw=black,line width=0.48pt] (15,-1) circle (0.3);
    \filldraw[fill=kgray,draw=black,line width=0.48pt] (16,-1) circle (0.3);
    \draw[black,line width=0.48pt] (15.788,-0.788) -- (16.212,-1.212) (15.788,-1.212) -- (16.212,-0.788);
    \filldraw[fill=white,draw=black,line width=0.48pt] (1,-2) circle (0.3);
    \filldraw[fill=white,draw=black,line width=0.48pt] (2,-2) circle (0.3);
    \filldraw[fill=kgray,draw=black,line width=0.48pt] (3,-2) circle (0.3);
    \draw[black,line width=0.48pt] (2.788,-1.788) -- (3.212,-2.212) (2.788,-2.212) -- (3.212,-1.788);
    \filldraw[fill=white,draw=black,line width=0.48pt] (4,-2) circle (0.3);
    \filldraw[fill=white,draw=black,line width=0.48pt] (5,-2) circle (0.3);
    \filldraw[fill=white,draw=black,line width=0.48pt] (6,-2) circle (0.3);
    \filldraw[fill=white,draw=black,line width=0.48pt] (7,-2) circle (0.3);
    \filldraw[fill=white,draw=black,line width=0.48pt] (8,-2) circle (0.3);
    \filldraw[fill=white,draw=black,line width=0.48pt] (9,-2) circle (0.3);
    \filldraw[fill=white,draw=black,line width=0.48pt] (10,-2) circle (0.3);
    \filldraw[fill=white,draw=black,line width=0.48pt] (11,-2) circle (0.3);
    \filldraw[fill=white,draw=black,line width=0.48pt] (12,-2) circle (0.3);
    \filldraw[fill=white,draw=black,line width=0.48pt] (13,-2) circle (0.3);
    \filldraw[fill=white,draw=black,line width=0.48pt] (14,-2) circle (0.3);
    \filldraw[fill=kgray,draw=black,line width=0.48pt] (15,-2) circle (0.3);
    \draw[black,line width=0.48pt] (14.788,-1.788) -- (15.212,-2.212) (14.788,-2.212) -- (15.212,-1.788);
    \filldraw[fill=kgray,draw=black,line width=0.48pt] (16,-2) circle (0.3);
    \draw[black,line width=0.48pt] (15.788,-1.788) -- (16.212,-2.212) (15.788,-2.212) -- (16.212,-1.788);
    \filldraw[fill=kgray,draw=black,line width=0.48pt] (1,-3) circle (0.3);
    \draw[black,line width=0.48pt] (0.788,-2.788) -- (1.212,-3.212) (0.788,-3.212) -- (1.212,-2.788);
    \filldraw[fill=white,draw=black,line width=0.48pt] (2,-3) circle (0.3);
    \filldraw[fill=white,draw=black,line width=0.48pt] (3,-3) circle (0.3);
    \filldraw[fill=white,draw=black,line width=0.48pt] (4,-3) circle (0.3);
    \filldraw[fill=white,draw=black,line width=0.48pt] (5,-3) circle (0.3);
    \filldraw[fill=white,draw=black,line width=0.48pt] (6,-3) circle (0.3);
    \filldraw[fill=white,draw=black,line width=0.48pt] (7,-3) circle (0.3);
    \filldraw[fill=white,draw=black,line width=0.48pt] (8,-3) circle (0.3);
    \filldraw[fill=white,draw=black,line width=0.48pt] (9,-3) circle (0.3);
    \filldraw[fill=white,draw=black,line width=0.48pt] (10,-3) circle (0.3);
    \filldraw[fill=white,draw=black,line width=0.48pt] (11,-3) circle (0.3);
    \filldraw[fill=white,draw=black,line width=0.48pt] (12,-3) circle (0.3);
    \filldraw[fill=white,draw=black,line width=0.48pt] (13,-3) circle (0.3);
    \filldraw[fill=kgray,draw=black,line width=0.48pt] (14,-3) circle (0.3);
    \draw[black,line width=0.48pt] (13.788,-2.788) -- (14.212,-3.212) (13.788,-3.212) -- (14.212,-2.788);
    \filldraw[fill=kgray,draw=black,line width=0.48pt] (15,-3) circle (0.3);
    \draw[black,line width=0.48pt] (14.788,-2.788) -- (15.212,-3.212) (14.788,-3.212) -- (15.212,-2.788);
    \filldraw[fill=white,draw=black,line width=0.48pt] (16,-3) circle (0.3);
    \filldraw[fill=white,draw=black,line width=0.48pt] (1,-4) circle (0.3);
    \filldraw[fill=kgray,draw=black,line width=0.48pt] (2,-4) circle (0.3);
    \draw[black,line width=0.48pt] (1.788,-3.788) -- (2.212,-4.212) (1.788,-4.212) -- (2.212,-3.788);
    \filldraw[fill=white,draw=black,line width=0.48pt] (3,-4) circle (0.3);
    \filldraw[fill=white,draw=black,line width=0.48pt] (4,-4) circle (0.3);
    \filldraw[fill=white,draw=black,line width=0.48pt] (5,-4) circle (0.3);
    \filldraw[fill=white,draw=black,line width=0.48pt] (6,-4) circle (0.3);
    \filldraw[fill=white,draw=black,line width=0.48pt] (7,-4) circle (0.3);
    \filldraw[fill=white,draw=black,line width=0.48pt] (8,-4) circle (0.3);
    \filldraw[fill=white,draw=black,line width=0.48pt] (9,-4) circle (0.3);
    \filldraw[fill=kgray,draw=black,line width=0.48pt] (10,-4) circle (0.3);
    \draw[black,line width=0.48pt] (9.788,-3.788) -- (10.212,-4.212) (9.788,-4.212) -- (10.212,-3.788);
    \filldraw[fill=kgray,draw=black,line width=0.48pt] (11,-4) circle (0.3);
    \draw[black,line width=0.48pt] (10.788,-3.788) -- (11.212,-4.212) (10.788,-4.212) -- (11.212,-3.788);
    \filldraw[fill=kgray,draw=black,line width=0.48pt] (12,-4) circle (0.3);
    \draw[black,line width=0.48pt] (11.788,-3.788) -- (12.212,-4.212) (11.788,-4.212) -- (12.212,-3.788);
    \filldraw[fill=kgray,draw=black,line width=0.48pt] (13,-4) circle (0.3);
    \draw[black,line width=0.48pt] (12.788,-3.788) -- (13.212,-4.212) (12.788,-4.212) -- (13.212,-3.788);
    \filldraw[fill=kgray,draw=black,line width=0.48pt] (14,-4) circle (0.3);
    \draw[black,line width=0.48pt] (13.788,-3.788) -- (14.212,-4.212) (13.788,-4.212) -- (14.212,-3.788);
    \filldraw[fill=white,draw=black,line width=0.48pt] (15,-4) circle (0.3);
    \filldraw[fill=white,draw=black,line width=0.48pt] (16,-4) circle (0.3);
    \filldraw[fill=white,draw=black,line width=0.48pt] (1,-5) circle (0.3);
    \filldraw[fill=white,draw=black,line width=0.48pt] (2,-5) circle (0.3);
    \filldraw[fill=kgray,draw=black,line width=0.48pt] (3,-5) circle (0.3);
    \draw[black,line width=0.48pt] (2.788,-4.788) -- (3.212,-5.212) (2.788,-5.212) -- (3.212,-4.788);
    \filldraw[fill=white,draw=black,line width=0.48pt] (4,-5) circle (0.3);
    \filldraw[fill=white,draw=black,line width=0.48pt] (5,-5) circle (0.3);
    \filldraw[fill=white,draw=black,line width=0.48pt] (6,-5) circle (0.3);
    \filldraw[fill=white,draw=black,line width=0.48pt] (7,-5) circle (0.3);
    \filldraw[fill=white,draw=black,line width=0.48pt] (8,-5) circle (0.3);
    \filldraw[fill=kgray,draw=black,line width=0.48pt] (9,-5) circle (0.3);
    \draw[black,line width=0.48pt] (8.788,-4.788) -- (9.212,-5.212) (8.788,-5.212) -- (9.212,-4.788);
    \filldraw[fill=kgray,draw=black,line width=0.48pt] (10,-5) circle (0.3);
    \draw[black,line width=0.48pt] (9.788,-4.788) -- (10.212,-5.212) (9.788,-5.212) -- (10.212,-4.788);
    \filldraw[fill=white,draw=black,line width=0.48pt] (11,-5) circle (0.3);
    \filldraw[fill=white,draw=black,line width=0.48pt] (12,-5) circle (0.3);
    \filldraw[fill=white,draw=black,line width=0.48pt] (13,-5) circle (0.3);
    \filldraw[fill=white,draw=black,line width=0.48pt] (14,-5) circle (0.3);
    \filldraw[fill=white,draw=black,line width=0.48pt] (15,-5) circle (0.3);
    \filldraw[fill=white,draw=black,line width=0.48pt] (16,-5) circle (0.3);
    \filldraw[fill=kgray,draw=black,line width=0.48pt] (1,-6) circle (0.3);
    \draw[black,line width=0.48pt] (0.788,-5.788) -- (1.212,-6.212) (0.788,-6.212) -- (1.212,-5.788);
    \filldraw[fill=white,draw=black,line width=0.48pt] (2,-6) circle (0.3);
    \filldraw[fill=white,draw=black,line width=0.48pt] (3,-6) circle (0.3);
    \filldraw[fill=kgray,draw=black,line width=0.48pt] (4,-6) circle (0.3);
    \draw[black,line width=0.48pt] (3.788,-5.788) -- (4.212,-6.212) (3.788,-6.212) -- (4.212,-5.788);
    \filldraw[fill=kgray,draw=black,line width=0.48pt] (5,-6) circle (0.3);
    \draw[black,line width=0.48pt] (4.788,-5.788) -- (5.212,-6.212) (4.788,-6.212) -- (5.212,-5.788);
    \filldraw[fill=kgray,draw=black,line width=0.48pt] (6,-6) circle (0.3);
    \draw[black,line width=0.48pt] (5.788,-5.788) -- (6.212,-6.212) (5.788,-6.212) -- (6.212,-5.788);
    \filldraw[fill=kgray,draw=black,line width=0.48pt] (7,-6) circle (0.3);
    \draw[black,line width=0.48pt] (6.788,-5.788) -- (7.212,-6.212) (6.788,-6.212) -- (7.212,-5.788);
    \filldraw[fill=kgray,draw=black,line width=0.48pt] (8,-6) circle (0.3);
    \draw[black,line width=0.48pt] (7.788,-5.788) -- (8.212,-6.212) (7.788,-6.212) -- (8.212,-5.788);
    \filldraw[fill=kgray,draw=black,line width=0.48pt] (9,-6) circle (0.3);
    \draw[black,line width=0.48pt] (8.788,-5.788) -- (9.212,-6.212) (8.788,-6.212) -- (9.212,-5.788);
    \filldraw[fill=white,draw=black,line width=0.48pt] (10,-6) circle (0.3);
    \filldraw[fill=white,draw=black,line width=0.48pt] (11,-6) circle (0.3);
    \filldraw[fill=white,draw=black,line width=0.48pt] (12,-6) circle (0.3);
    \filldraw[fill=white,draw=black,line width=0.48pt] (13,-6) circle (0.3);
    \filldraw[fill=white,draw=black,line width=0.48pt] (14,-6) circle (0.3);
    \filldraw[fill=white,draw=black,line width=0.48pt] (15,-6) circle (0.3);
    \filldraw[fill=white,draw=black,line width=0.48pt] (16,-6) circle (0.3);
    \node[lab] at (1,-0.45) {1};
    \node[lab] at (2,-0.45) {2};
    \node[lab] at (3,-0.45) {3};
    \node[lab] at (4,-0.45) {4};
    \node[lab] at (5,-0.45) {5};
    \node[lab] at (6,-0.45) {6};
    \node[lab] at (7,-0.45) {7};
    \node[lab] at (8,-0.45) {8};
    \node[lab] at (9,-0.45) {9};
    \node[lab] at (10,-0.45) {10};
    \node[lab] at (11,-0.45) {11};
    \node[lab] at (12,-0.45) {12};
    \node[lab] at (13,-0.45) {13};
    \node[lab] at (14,-0.45) {14};
    \node[lab] at (15,-0.45) {15};
    \node[lab] at (16,-0.45) {16};
    \node[lab] at (0.45,-1) {1};
    \node[lab] at (0.45,-2) {2};
    \node[lab] at (0.45,-3) {3};
    \node[lab] at (0.45,-4) {4};
    \node[lab] at (0.45,-5) {5};
    \node[lab] at (0.45,-6) {6};
    \node[capt,anchor=north] at (8.50,-6.45) {(b)};
  \end{scope}
\end{tikzpicture}

\caption{The $5\times13$ test board of $S_8$ in (a) and the blackout $S_{16}$ of the $6\times16$ grid in (b). Gray crossed circles are blacked out and white circles are visible, and rows are numbered from the top.}
\label{fig:boards}
\end{figure}

\begin{theorem}[five rows, (a) and (c) computer-assisted]\label{thm:fiverow}
\begin{enumerate}[(a)]
\item We have $f(5,m)=m+5$ for $m=5$, $6$, and $7$.
\item We have $f(5,m)=m+6$ for every $m\ge8$, attained by $S_8$ followed by $m-8$ columns black in
row $1$ only.
\item The $5\times8$ grid has exactly four maximum blackouts, the images of $S_8$ under the row
reflection and the column reversal.
\end{enumerate}
\end{theorem}

\begin{proof}
(b) The upper bound is Theorem~\ref{thm:rowbound}. Apply Theorem~\ref{thm:collar} to $S_8$
with the padding row $1$ on the right. Since $W_5=5$, the test board is the $5\times13$
board of Figure~\ref{fig:boards}(a). Three independent enumerators find no repeated
presentation among its $954$ rectangles, and Lean proves part~(b) in full
(Section~\ref{sec:verify}).

(a) The lower bounds are $S_0$ of Theorem~\ref{thm:fivefive} below for $m=5$ and stored
blackouts of $h+m$ points for $m=6,7$. For $m=5$ the upper bound is
Corollary~\ref{cor:small-m}. For $m=6,7$ a valid blackout of $m+6$ points would be one of the
$10\,(m)_6\,5^{m-6}$ candidates of Theorem~\ref{thm:equality}, and SAT encodings exclude all of
them, with checked certificates and independent cross-checks (Section~\ref{sec:verify}).

(c) Among the $10\,(8)_6\,5^2=5{,}040{,}000$ candidates of Theorem~\ref{thm:equality}, two
SAT enumerations and a column search find exactly the four images of $S_8$.
\end{proof}
\subsection{Six rows and small grids}

Let $S_{16}$ be the blackout of the $6\times16$ grid in Figure~\ref{fig:boards}(b). Its doubles, in the columns $1,2,3,6,9,10,14,15,16$, carry the edges of $K_{3,3}$ between $\{1,3,5\}$ and $\{2,4,6\}$.

\begin{theorem}[six rows, (a) computer-assisted]\label{thm:sixrow}
\begin{enumerate}[(a)]
\item For $8\le m\le14$ we have $f(6,m)\le m+8$.
\item For every $m\ge15$ we have $f(6,m)=m+9$.
\end{enumerate}
\end{theorem}

\begin{proof}
(a) For $m=8$ this is Corollary~\ref{cor:small-m}. For $9\le m\le14$, by
Theorem~\ref{thm:equality} the doubles of a valid blackout of $m+9$ points carry the edges of
$K_{3,3}$ for one of ten row partitions. Two SAT encoders find none for any width and
partition, with checked LRAT certificates and a third check at every width
(Section~\ref{sec:verify}).

(b) The upper bound is Theorem~\ref{thm:rowbound}. A valid blackout of $24$ points on the
$6\times15$ grid is stored in the repository~\cite{repo}. For $m\ge16$ apply
Theorem~\ref{thm:collar} to $S_{16}$ with the padding row $1$ on the right. Since $W_6=7$,
the test board adds $(17,1),\dots,(23,1)$, and three enumerators find no repeated
presentation among its $4{,}627$ rectangles. Lean proves part~(b) in full.
\end{proof}

% From sections/sec07_five.tex, round 3 cut 2026-10-19 (hand-check sentence before the smallest-area remark).
Checking the $364$ and $592$ rectangles of $S_8$ and of this blackout shows that both are
valid, so the guess of Section~\ref{sec:intro} fails without a solver.
